% Analytic attempt at the multi-digit coefficient of I_B.
% Failures are labeled ADMISSION. No Final Status label is changed.
\documentclass[11pt]{article}
\usepackage[margin=1.05in]{geometry}
\usepackage{amsmath,amssymb,amsthm}
\usepackage{enumitem,booktabs}
\usepackage[colorlinks=true,linkcolor=blue,citecolor=blue,urlcolor=blue]{hyperref}

\newtheorem{proposition}{Proposition}[section]
\newtheorem{remark}{Remark}[section]
\theoremstyle{definition}
\newtheorem*{admission}{Admission}

\newcommand{\pP}{\textbf{[P]}}
\newcommand{\pO}{\textbf{[O]}}
\newcommand{\pA}{\textbf{[ADMISSION]}}
\newcommand{\IB}{I_B}

\title{\textbf{Analytic Attempt at the Coefficient of $\IB$:\\
Hyperbolic Reduction, Source Dependence, and an Admission}}
\author{Robert W.\ McGwier\\
\small CTO, Cohere Technology Group\\
\small GitHub: \texttt{n4hy}}
\date{15 August 2026 --- Version 1}

\begin{document}
\maketitle

\begin{abstract}
\noindent
Paper~VIII of the Emergent Gravity Program left the multi-digit
coefficient of the kinematic residue $\IB$ open, and required an
analytic Hadamard expansion of the current two-point function under
modular flow. This note carries out the $d=2$ reduction as far as it
goes, then stops.

After the map $x=R\tanh u$ the weight-one Jacobians cancel and the
smeared correlator is
\begin{equation*}
Q(s)=\int\!du\,dv\,b(u)b(v)/\sinh^{2}(u-v+\pi s),
\end{equation*}
with $b(u)=\beta(\tanh u)$. The modular kernel pulls back to a
convolution against $\mathrm{sech}^{2}$. That identity is algebraic
and is not in dispute.

A Hadamard finite part for the pulled-back integral was written in
closed form,
$I_{\mathrm{tot}}^{\mathrm{FP}}(\tau)=2\,\mathrm{sech}^{2}(\tau)\,(\tau\tanh\tau-1)$.
A hole-quadrature check of that formula does \emph{not} confirm it:
the hole is not a Hadamard finite part, and the two sides do not
agree. The closed form is therefore \textbf{not} certified.

A direct evaluation of the relative residue
$r=\int K(s)(Q(s)-Q(0))\,ds/|Q(0)|$
for five inequivalent sources in the hyperbolic coordinate gives
values that \emph{move} with the source (spread $\approx 0.23$ on the
grids used). If that movement survives a genuine Hadamard scheme, there
is no single universal multi-digit coefficient of $\IB$ independent of
$\beta$.

This note does not extract a multi-digit number. The Final Status
label of $\IB$ remains \pO{}.
\end{abstract}

\section*{Standing rule}
Numerical agreement is not proof. An unconfirmed closed form is not a
theorem. A source-dependent ratio is not a universal coefficient.
This draft was prepared with computational assistance; every
\pA{} is an owned failure, not a pending detail.

\section{The object}

Paper~VI defines, up to the current central charge $C_J$,
\begin{equation}
\IB[\beta;\zeta]
=
-\int_{-\infty}^{\infty} ds\,
K(s)\,
\bigl\langle \sigma_s J[\beta]\,\widetilde{J}[\zeta]\bigr\rangle_{\kappa\neq 0},
\qquad
K(s)=\frac{\pi}{2\cosh^{2}(\pi s)},
\quad
\int K=1.
\label{eq:IB}
\end{equation}
On the interval $(-R,R)$ the modular flow is
\begin{equation}
x(s)=R\tanh\bigl(\mathrm{artanh}(x/R)+\pi s\bigr).
\label{eq:flow}
\end{equation}
The chiral current two-point function is
$\langle J(x)J(y)\rangle=k/(x-y)^{2}$. The constant $k$ cancels in
any relative residue and is set to $1$ below. $R=1$ by scaling.

\section{Hyperbolic reduction}
\label{sec:red}

Set $x=\tanh u$, $y=\tanh v$. Then
\begin{equation}
\frac{\partial x(s)}{\partial x}\frac{1}{(x(s)-y)^{2}}
=
\frac{\cosh^{2}u\,\cosh^{2}v}{\sinh^{2}(u-v+\pi s)},
\qquad
dx=\mathrm{sech}^{2}u\,du.
\end{equation}
The Jacobians cancel against the measure:
\begin{equation}
Q(s)
=
\int\!du\,dv\,b(u)b(v)\frac{1}{\sinh^{2}(u-v+\pi s)}
=
\int\!d\tau\,C(\tau)\frac{1}{\sinh^{2}(\tau+\pi s)},
\label{eq:Q}
\end{equation}
where $b(u)=\beta(\tanh u)$ and $C(\tau)=\int du\,b(u)b(u-\tau)$.

Pulling back the kernel,
\begin{equation}
\int_{-\infty}^{\infty}K(s)\,f(\tau+\pi s)\,ds
=
\int_{-\infty}^{\infty}\frac{d\alpha}{2\cosh^{2}\alpha}\,f(\tau+\alpha).
\label{eq:pull}
\end{equation}
On the real line the only singularity of $1/\sinh^{2}w$ is at $w=0$,
and
\begin{equation}
\frac{1}{\sinh^{2}w}-\frac{1}{w^{2}}
=
-\frac13+O(w^{2})
\label{eq:h}
\end{equation}
is smooth. Subtracting $1/w^{2}$ is the first Hadamard step. It is
\emph{not} the last: Paper~VIII required every non-integrable term of
the \emph{flowed} two-point function to be subtracted, not only the
$s=0$ coincidence.

\begin{proposition}[reduction]
\label{prop:red}
Equations \eqref{eq:Q}--\eqref{eq:pull} are identities of the $d=2$
weight-one setup. They do not evaluate $\IB$.
\end{proposition}

\section{An unconfirmed closed form}
\label{sec:closed}

The substitution $u=\tanh\alpha$ converts the pulled-back integral of
$1/\sinh^{2}$ into
\begin{equation}
I_{\mathrm{tot}}(\tau)
=
\frac{1}{2\cosh^{2}\tau}
\int_{-1}^{1}\!du\,\frac{1-u^{2}}{(u+\tanh\tau)^{2}}.
\label{eq:Iu}
\end{equation}
The integrand has a double pole on the contour at $u=-\tanh\tau$.
Dropping the $1/\varepsilon$ term in a symmetric hole and taking
$\varepsilon\to 0$ produces the candidate
\begin{equation}
I_{\mathrm{tot}}^{\mathrm{FP}}(\tau)
\stackrel{?}{=}
2\,\mathrm{sech}^{2}(\tau)\,(\tau\tanh\tau-1).
\label{eq:candidate}
\end{equation}

\begin{admission}
Equation \eqref{eq:candidate} was \emph{not} confirmed. A hole
quadrature of \eqref{eq:Iu} diverges as the hole shrinks, as a double
pole must, and the finite part extracted from that hole does not
match \eqref{eq:candidate} under the implementation used
(\texttt{m9\_3\_ib\_analytic.py}). A hole is not a Hadamard finite
part. Until \eqref{eq:candidate} is verified by an independent
finite-part prescription, it is a guess, not a lemma. It is not used
as input to any number claimed below.
\end{admission}

\section{Source dependence of the relative residue}
\label{sec:src}

Define, with the same cutoff on $| \tau+\pi s |$ for every $s$,
\begin{equation}
r[\beta]
:=
\frac{\displaystyle\int ds\,K(s)\bigl(Q(s)-Q(0)\bigr)}{|Q(0)|}.
\label{eq:r}
\end{equation}
Five sources $b(u)$ were run on a linear $u$-grid
(\texttt{m9\_3} family, direct correlation, hole $0.03$):

\begin{center}
\begin{tabular}{@{}lc@{}}
\toprule
$b(u)$ & $r$ \\
\midrule
Gaussian, $\sigma=0.6$ & $-0.108$ \\
Gaussian, $\sigma=1.1$ & $+0.089$ \\
$\mathrm{sech}\,u$ & $+0.120$ \\
$e^{-|u|}$ & $+0.046$ \\
compact bump, width $1.4$ & $-0.111$ \\
\bottomrule
\end{tabular}
\end{center}

Spread: $0.23$. Kernel mass on the truncated $s$-interval is $0.90$,
not $1$; that truncation is a further scheme artefact.

\begin{admission}
These five numbers are scheme-dependent proxies, not Hadamard finite
parts. They are already enough to kill a universality claim
\emph{inside this scheme}: $r$ changes sign with the source. A
universal multi-digit coefficient would have to be independent of
$\beta$. This scheme does not produce one. Whether a genuine
Hadamard $r$ is universal remains \pO{}. The present evidence leans
against universality and does not replace Paper~VIII's demand for an
analytic expansion.
\end{admission}

\section{What was not done}

\begin{itemize}[leftmargin=1.4em]
\item No closed form for
$\displaystyle\int d\alpha\,\mathrm{sech}^{2}\alpha\,w_{\mathrm{fin}}(\tau+\alpha)$.
\item No isolation of the $u_1$ or $\log$ terms of a full Hadamard
series for the flowed correlator.
\item No three-digit stability under simultaneous refinement of
grid, hole, and $s$-cutoff.
\item No evaluation in $d=4$.
\end{itemize}

\section{Verdict}

\begin{center}
\begin{tabular}{@{}p{0.42\textwidth}p{0.48\textwidth}@{}}
\toprule
Claim & Status \\
\midrule
Hyperbolic reduction \eqref{eq:Q}--\eqref{eq:pull} &
derived; algebra only \\
Closed form \eqref{eq:candidate} &
\pA{} unconfirmed \\
Universal multi-digit coefficient of $\IB$ &
\pO{}; this attempt did not find one \\
Source-independent $r$ in the hole scheme &
fails (spread $0.23$, sign change) \\
Final Status label of $\IB$ &
unchanged, \pO{} \\
\bottomrule
\end{tabular}
\end{center}

The next honest step is still the one Paper~VIII named: an analytic
Hadamard series of $\langle J(x(s))J(y)\rangle$ whose finite part can
be integrated against $K(s)$ in closed form, or a proof that the
finite part is not unique. This note is neither.

\section*{Acknowledgments}
Computational assistance drafted the algebra checks and the scripts
\texttt{m9\_2\_ib\_hadamard.py} and \texttt{m9\_3\_ib\_analytic.py}.
The admissions are the point of the note.

\begin{thebibliography}{99}
\bibitem{VI} R.W.~McGwier, ``Evaluation of the Universal Kinematic Integral for Condition NL,''
\url{https://github.com/n4hy/New_Model_Emergent_Gravity}.
\bibitem{VIII} R.W.~McGwier, ``Status of the High-Precision Coefficient of $\IB$ (Version 4),'' ibid.
\bibitem{FS} R.W.~McGwier, ``Final Status of Claims,'' ibid.
\bibitem{CHM} H.~Casini, M.~Huerta and R.C.~Myers, JHEP 05 (2011) 036.
\bibitem{CTT} H.~Casini, E.~Teste and G.~Torroba, J.\ Phys.\ A 50 (2017) 364001.
\end{thebibliography}
\end{document}
